\documentclass[11pt]{article}

% ----- Layout -----
\usepackage[margin=1in]{geometry}
\usepackage{microtype}
\usepackage{setspace}
\setstretch{1.15}

% ----- Math -----
\usepackage{amsmath, amssymb, amsthm, mathtools}

% ----- Figures/tables -----
\usepackage{graphicx}
\usepackage{booktabs}
\usepackage{caption}
\usepackage{subcaption}

% ----- Links + smart refs -----
\usepackage[hidelinks]{hyperref}

% ----- BibLaTeX (biber) -----
\usepackage[
  backend=biber,
  style=numeric-comp,
  maxcitenames=2,
  maxbibnames=99,
  doi=true,
  isbn=false,
  url=true
]{biblatex}
\addbibresource{references.bib}

% ----- Title info -----
\title{How Model Structure and Parameter Choices Influence Tumor Immune Dynamics and Treatment Response}
\author{AMATH 881 \\Rosie DeFazio\\ 20628115}
\date{2026-04-13}

% ----- Optional: theorem envs -----
\newtheorem{theorem}{Theorem}
\newtheorem{proposition}{Proposition}
\theoremstyle{definition}
\newtheorem{definition}{Definition}

\begin{document}
\maketitle

%\begin{abstract}
%A 5--10 line summary of your motivation, approach, and key results.
%Cite related work like \textcite{knuth1984} or parenthetical \parencite{lamport1994}.
%\end{abstract}

\section{Introduction}
From the first report in 1904 of tumor regression during viral infections to various immune cells being discovered and clinical interventions being approved; the role and interactions of the immune system in regards to cancer 
has been well researched thoroughly scientific methods \parencite{aacr2025history}, revealing complex interactions. This paper builds on the model by Hernandez-Lopez et al. \parencite{ERX2026} inspired by Elewaut et al. \parencite{ERX2026} and supplemented by Reis et al. \parencite{reis2021bringing} further extending by exploring the parameter choices described in the article and extending 
the model to include macrophage and chemotherapy interactions; allowing for further biological interpretation of the complex system. In Section~2, the original model from \parencite{ERX2026} is revisited, highlighting incompatible parameter choices with
interpretation along side investigation of vague parameter definition. Section~3 explores a macrophage extension; identifying the tumor-free equilibrium and characterizing its stability.
A chemotherapy extension is presented in Section~4, also contained exploration of the tumor-free equilibrium point. 
\section{Original Model}
Below is the model introduced in Hernandez-Lopez et al. \parencite{ERX2026} with $C(t)$ representing the density of cancer cells, $T(t)$ the density of T-Lymphocytes cells, $M_n(t)$ the density of monocytes and 
$M(t)$ the density of macrophages. This work was inspired by the experiments preformed in Elewaut et al. \parencite{El2024} which highlighted the importance of the MCHI cross dressing between monocytes and T-cells 
captured by the term $\beta MT$.
\\
$$
\begin{aligned}
\frac{dC}{dt} &= rC (1 - bC) (C - \gamma) - \alpha CT, \\
\frac{dT}{dt} &= \lambda C + \beta MT - \delta T, \\
\frac{dM}{dt} &= \rho - \eta M - \omega CM.
\end{aligned}
$$
\\
The experiments suggested that by the MCHI cross dressing mechanism, monocytes can activate the stimulation of T-cells. Further research from DeNardo et al. \parencite{DeNardo2019}
suggest that monocytes can also suppress the stimulation of T-cells thus allowing the parameter $\beta$ to take on negative values. This has an effect on the tumor free equilibrium state for this model as stated in \parencite{ERX2026}: 
the tumor-free equilibrium $E_0 = (0, 0, \rho/\eta)$. The stability of this equilibrium is as follows.
\begin{enumerate}
  \item[\textbf{(a)}] $E_0$ is locally asymptotically stable if $\beta\rho < \delta\eta$.
  \item[\textbf{(b)}] $E_0$ is a saddle if $\beta\rho > \delta\eta$.
  \item[\textbf{(c)}] $E_0$ is a degenerate equilibrium if $\beta\rho = \delta\eta$.
\end{enumerate}
Thus contributing to the following proposition: 
\begin{proposition}
$E_0$ is locally asymptotically stable when $\beta<0$ and $\rho, \eta, \delta >0$
\end{proposition}
\begin{proof}
Using proposition 2.1 from \parencite{ERX2026}, when $\beta < 0 \implies \beta\rho < \delta\eta$ for strictly positive values of $\rho, \eta, \delta$.
\end{proof}

\noindent Based on the literature from \parencite{DeNardo2019}, this seems to be a biologically plausible state, yet no dynamics of negative $\beta$ values where explored in the original paper. 
Another note on the model, is that within the paper, it was stated that $\eta$ can be set to zero, justifying that the $-\omega CM$ term would capture the apoptosis mechanism captured by the $- \eta M$ term due to the dominant nature of cancer mediated monocyte death. 
Yet, upon further inspection, setting $\eta= 0$ causes $\rho = 0$ thus biologically invalidating the model as this implied that monocytes have no production mechanism.    

\section{Macrophage Extension}
The system for macrophage extension as follows: 
\\
$$
\begin{aligned}
\frac{dC}{dt} &= rC(1-bC)(C-\gamma) - \alpha CT - aCM_c \\
\frac{dT}{dt} &= \lambda C + \beta M_n T - \delta T + \phi T M_c \\
\frac{dM_n}{dt} &= \rho - \eta M_n - \omega CM_n - gM_nM_c \\
\frac{dM_c}{dt} &=  -d M_cM_n - eTM_c + h 
\end{aligned}
$$
\\
This is the system for macrophage extension. $C(t)$ denotes the density of cancer cells, $T(t)$ the density of T-Lymphocytes cells, $M_n(t)$ the density of monocytes and 
$M_c(t)$ the density of macrophages. The base of the model comes form Hernandez-Lopez, as described in section 2. 
The additional term in the cancer cell equation of $-aCM_c$ captures the 
inhibitory effects of macrophages on the cancer cells \parencite{reis2021bringing}.  $\phi T M_c$ represents the stimulatory effect of macrophages on T Lymphocytes.
Whereas, the term $-gM_nM_c$ models the maturation of monocytes into macrophages. The additional equation $\frac{dM_c}{dt}$ aims to capture simplified 
dynamics of the macrophage interactions within the tumor microenvironment with $-d M_cM_n$ representing the reduction in monocytes after expansion of inflammatory monocytes as noted in Elewaut et al., $h$ refers to a 
general recruitment rate for macrophages and $-eTM_c$ capturing the observed absence of macrophages in tumor areas with T cells absent as well \parencite{El2024}. 
\\
\begin{proposition}
  The point \(P_0 = \left( 0, 0, M_n^* =\frac{d\rho - gh}{dn}, M_c^* =\frac{hn}{d\rho - gh} \right)\) is the tumor free equilibrium of the system with 
  stability as follows: 
  \begin{enumerate}
    \item $P_0$ is asymptotically stable if \(\beta M_n^* + \phi M_c^* < \delta\)
    \item $P_0$ is saddle if \(\beta M_n^* + \phi M_c^* > \delta\)
    \item $P_0$ is degenerate if \(\beta M_n^* + \phi M_c^* = \delta\)
  \end{enumerate}
\end{proposition}

\begin{proof}
To analyze the stability of $P_0$ it is necessary to compute the Jacobian of the system: 

\begin{center}
$$J = \left(\begin{smallmatrix}
-r(\gamma + C(-2 + 3bC - 2b\gamma)) - \alpha T- aM_c & -\alpha C & 0 & -aC \\
\lambda & \beta M_n - \delta + \phi M_c & \beta T & \phi T \\
-\omega M_n & 0 & -\eta - \omega C - gM_c & -gM_n \\
0 & -eM_c & -dM_c & -dM_n - eT
\end{smallmatrix}\right)$$
\end{center}
Substituting the tumor free equilibrium gives: 
\\
\begin{center}
  $$ \begin{pmatrix}
-r\gamma - aM_c^* & 0 & 0 & 0 \\
\lambda & \beta M_n^* - \delta + \phi M_c^* & 0 & 0 \\
-\omega M_n^* & 0 & -n - gM_c^* & -gM_n^* \\
0 & -eM_c^* & -dM_c^* & -dM_n^*
\end{pmatrix}
  $$
\end{center}

\noindent The eigenvalues of this matrix are: 
$$
\begin{aligned}
  \xi_1 &= -r\gamma - aM_c^*\\
  \xi_2 &= \beta M_n^* - \delta + \phi M_c^*\\
  \xi_{3,4} &= \frac{-(n + gM_c^* + dM_n^*) \pm \sqrt{(n + gM_c^* + dM_n^*)^2 - 4ndM_n^*}}{2}
\end{aligned}
$$

\noindent For stability, all eigenvalues must have negative real parts. For $\xi_{1,3,4}$ this is true for all values of the parameter, leaving 
$\xi_2$ to be the predictor of stability in the system. 
\end{proof}

\section{Chemotherapy Extension}
Below is the model which further extends the model from Hernandez-Lopez \parencite{ERX2026} by introducing Chemotherapy terms alongside the macrophage extension discussed in section 3.
\\
$$
\begin{aligned}
\frac{dC}{dt} &= rC(1 - bC)(C - \gamma) -  \alpha CT - \kappa CD - aCM_c  \\
\frac{dT}{dt} &= \lambda C + \beta M_nT - \delta T + \phi TM_c - qTD\\
\frac{dM_n}{dt} &= p - \eta M_n - \omega C M_n - g M_n M_c \\
\frac{dM_c}{dt} &= -d M_c M_n - e T M_c + h - i M_c D\\
\frac{dD}{dt} &= u - \tau D
\end{aligned}
$$
\\
$C(t)$ denotes the density of cancer cells, $T(t)$ the density of T-Lymphocytes cells, $M_n(t)$ the density of monocytes,
$M_c(t)$ the density of macrophages and $D(t)$ represents the drug concentration within the tumor microenvironment over time. $\frac{dD}{dt}$ follows a standard 
model for drug concentration with the parameter $u$ capturing the peak drug concentration and $-\tau D$ modeling exponential decay.
Additional terms in the cancer, T-lymphocytes and macrophage equations were added to capture the effect of chemotherapy reducing healthy cells in addition to cancer cells \parencite{Sharma2024}. 
The monocyte equation does not have an additional term based  on the assumption chemotherapy does not have an effect on monocyte production.  
\\
\\
To assess the tumor free equilibrium it is necessary to split this paper moves forward with the assumption that the chemotherapy steady state also reaches eradication. This is biologically plausible as it 
represents the withdraw of chemotherapy after cancer eradication. With this assumption two cases unfold. The first is no T-cells are present within the tumor microenvironment, the second represents an active immune system 
in the TME after cancer eradication corresponding to $T^* >0$. Both cases are explored as follows: 

%--- T CELL NO EXIST EQUILIBRUM ---% 
\begin{proposition}
  The cancer, T-cell and drug free equilibrium point $P_0 =( 0, 0, \frac{dp - hg}{d\eta}, \frac{h\eta}{dp - hg}, 0)$ has stability characterization as follow: 
  \begin{enumerate}
    \item $P_0$ is asymptotically stable if \(\beta M_n^* + \phi M_c^* < \delta\)
    \item $P_0$ is saddle if \(\beta M_n^* + \phi M_c^* > \delta\)
    \item $P_0$ is degenerate if \(\beta M_n^* + \phi M_c^* = \delta\)
  \end{enumerate}
\end{proposition}
\begin{proof}

Finding the cancer and drug free equilibrium is found by setting $C^*, D^*, u, T^* =0$
\begin{align*}
\frac{dM_n}{dt} &= p - M_n(\eta + gM_c) = 0 \implies M_n^* = \frac{p}{\eta + gM_c^*}\\
\frac{dM_c}{dt} &= h - M_c(dM_n ) = 0 \implies M_c^* = \frac{h}{dM_n^*}\\
\end{align*}

\noindent The Jacobian of the chemotherapy system is: 
\begin{align*}
\left( \begin{smallmatrix}
-r(\gamma + C(-2 + 3bC - 2b\gamma)) - aM_c- \kappa D & -aC & 0 & 0 & -\kappa C \\
\lambda &  \beta M - \delta + \phi M_c - qD & 0 & \phi T & -qT \\
-\omega M_n & 0 & -\eta - \omega C - g M_c & -g M_n & 0 \\
0 & -e M_c & -d M_c & -d M_n - e T - i D & -i M_c \\
0 & 0 & 0 & 0 & -\tau
\end{smallmatrix} \right)
\end{align*}

\noindent Substituting the equilibrium point gives: 
\begin{align*}\left( \begin{smallmatrix}
-r\gamma - aM_c^* & 0 & 0 & 0 & 0 \\
\lambda & \beta M - \delta + \phi M_c^* & 0 & 0 & 0 \\
-\omega M_n^* & 0 & -(\eta + g M_c^*) & -g M_n^* & 0 \\
0 & -e M_c^* & -d M_c^* & -d M_n^* & -i M_c^* \\
0 & 0 & 0 & 0 & -\tau
\end{smallmatrix} \right) \end{align*}

\noindent The eigenvalues of this matrix can be found for $\xi_1, \xi_2, \xi_3$ since the form is in a upper diagonal matrix:
\begin{align*}
\xi_1 &= -r\gamma - aM_c^*\\
\xi_2 &= \beta M - \delta + \phi M_c^*\\
\xi_5 &= -\tau\\
\end{align*}

\noindent The eigenvalues for $\xi_3, \xi_4$ can be found by calculating the trace and determinant for the sub-matrix:

\begin{align*}
A_{sub} = \begin{pmatrix}
-(\eta + g M_c^*) & -g M_n^* \\
-d M_c^* & -d M_n^*
\end{pmatrix}
\end{align*}

\noindent and calculating the characteristic equation:
\begin{align*}
  \xi^2 - \text{Tr}(A_{sub})\xi + \det(A_{sub}) = 0
\end{align*}

\noindent Where:
\begin{align*}
\text{Tr}(A_{sub}) &= -(\eta + g M_c^* + d M_n^*)\\
\det(A_{sub}) &= (-\eta - g M_c^*)(-d M_n^*) - (-g M_n^*)(-d M_c^*) = d\eta M_n^*
\end{align*}
\noindent Using the quadratic formula:
\begin{align*}
\xi_{3,4} = \frac{-(\eta + g M_c^* + d M_n^*) \pm \sqrt{(\eta + g M_c^* + d M_n^*)^2 - 4d\eta M_n^*}}{2}
\end{align*}

\noindent Thus the equilibrium points are: 
\begin{align*}
\xi_1 &= -r\gamma - aM_c^*\\
\xi_2 &= \beta M - \delta + \phi M_c^*\\
\xi_3 &= \frac{-(\eta + g M_c^* + d M_n^*) + \sqrt{(\eta + g M_c^* + d M_n^*)^2 - 4d\eta M_n^*}}{2}\\
\xi_4 &= \frac{-(\eta + g M_c^* + d M_n^*) - \sqrt{(\eta + g M_c^* + d M_n^*)^2 - 4d\eta M_n^*}}{2}\\
\xi_5 &= -\tau
\end{align*}
\noindent which all eigenvalues except $\xi_2$ have negative real part thus creating the conditions for stability. 
\end{proof}

%---- TCELL POSTIVIE EQUILIBRUM -----%
\begin{proposition}
  The cancer, and drug free equilibrium point $P_0 = (0, \beta M_n^* + \phi M_c^* - \delta,M_n^* = \frac{p}{\eta + gM_c^*},M_c^* = \frac{h}{dM_n^* + eT^*}, 0) $ is stable when $a_1a_2 > a_3$
  \\
\\
  with
  \\ $a_1 =\eta + gM_c^* + dM_n^* + eT^*$,
  \\
  $a_2 = e\phi T^* M_c^*(\eta + gM_c^*)$, and 
  \\$a_3 = \eta d M_n^* + eT^*(\eta + gM_c^*) + e\phi T^* M_c^*$
\end{proposition}
\begin{proof}

Finding the cancer and drug free equilibrium is found by setting $C^*, D^*, u = 0$
\begin{align*}
\frac{dM_n}{dt} &= p - M_n(\eta + gM_c) = 0 \implies M_n^* = \frac{p}{\eta + gM_c^*}\\
\frac{dM_c}{dt} &= h - M_c(dM_n + eT) = 0 \implies M_c^* = \frac{h}{dM_n^* + eT^*}\\
\frac{dT}{dt} &= T(\beta M_n^* + \phi M_c^* - \delta) = 0 \implies T^* = \beta M_n^* + \phi M_c^* - \delta\\
\end{align*}

\noindent Using the Jacobian found in Proposition 3, and substituting in the equilibrium point gives: 
\begin{align*}
\begin{pmatrix}
-r\gamma - aM_c^* & 0 & 0 & 0 & 0 \\[4pt]
\lambda & 0 & 0 & \phi T^* & -qT^* \\[4pt]
-\omega M_n^* & 0 & -(\eta + gM_c^*) & -gM_n^* & 0 \\[4pt]
0 & -eM_c^* & -dM_c^* & -(dM_n^* + eT^*) & -iM_c^* \\[4pt]
0 & 0 & 0 & 0 & -\tau
\end{pmatrix}
\end{align*}
\noindent Upon inspection, eigenvalues $\xi_1= -r\gamma - aM_c^*, \xi_5 =-\tau$ can be obtained directly from the matrix as the other column entries are zero. The values will always be negative. 
The other eigenvalues can be obtained by inspecting the sub matrix and applying Routh-Hurwitz criteria as described in Appendix B of \parencite{murray2002mathematical}. 
\begin{align*}
\begin{pmatrix}
  0 & 0 & \phi T^* \\
0 & -(\eta + gM_c^*) & -gM_n^* \\
-eM_c^* & -dM_c^* & -(dM_n^* + eT^*)
\end{pmatrix}
\end{align*}
\noindent resulting in the sub characteristic equation of the form

\begin{align*}
  \lambda^3
+ a_1 \lambda^2
+ a_2\lambda
+ a_3
= 0
\end{align*}

\noindent where $a_1 =\eta + gM_c^* + dM_n^* + eT^*$ is the trace of the sub-matrix, $a_2 = e\phi T^* M_c^*(\eta + gM_c^*)$ is the determinant of the sub-matrix and $ a_3 = \eta d M_n^* + eT^*(\eta + gM_c^*) + e\phi T^* M_c^*$ is the sum of the 2x2 principal minors, alongside the fact that all parameters other than $\beta$ must be positive gives the criteria as needed.
\end{proof}

\section{Discussion}
\subsection*{Original Model}
A critical observation from Section~2 concerns the sign of the cross-dressing
parameter $\beta$. The original paper treated $\beta$ only produced bifurcation diagrams and hypersurfaces
for positive $\beta$
yet evidence from DeNardo et al.\ \parencite{DeNardo2019} 
supports a suppressive role for monocytes in certain tumor micro-environments,
thus allowing for negative values of $\beta$. Further evidence from the inspiration paper Elewaut et al. indicates that there could 
be no stimulatory effect indicating $\beta =0$. Proposition~1 formalities this: whenever
$\beta < 0$ and all remaining parameters are strictly positive, the condition
$\beta\rho < \delta\eta$ is automatically satisfied and the tumor-free
equilibrium $E_0$ is locally asymptotically stable. These parameter values were not
explored numerically in \parencite{ERX2026}. 
\\
\\
A second concern is the suggestion in \parencite{ERX2026} that $\eta$ may be
set to zero. As noted in Section~2, this forces $\rho = 0$ at the tumor-free
equilibrium, eliminating monocyte production entirely. Because monocytes are
continuously recruited from bone-marrow precursors under homeostatic conditions
\parencite{reis2021bringing}, a model with $\rho = 0$ lacks biological
grounding. A solution is to retain $\eta > 0$ as a small but non-zeroclearance rate, which preserves the monocyte source term while still
allowing cancer-mediated death ($-\omega CM$) to dominate during active tumor
growth.

\subsection*{Macrophage Extension}

The extension in Section~3 separates monocytes ($M_n$) from macrophages
($M_c$), reflecting monocyte-to-macrophage maturation \parencite{reis2021bringing}. The tumor-free equilibrium
$P_0 = (0,0,M_n^*,M_c^*)$ retains a clean stability criterion: the point is
locally asymptotically stable when $\beta M_n^* + \phi M_c^* < \delta$. This
threshold has an interpretation; stability requires that the combined
stimulatory effect of monocytes and macrophages on T-lymphocytes be outweighed
by T-cell death. A limitation to this is the model does not 
include the suppressive effects on macrophages on T-cells as defined in Reis et al. \parencite{reis2021bringing}. Future 
exploration can include further compartmentalization to describe M1 and M2 macrophages. 

\subsection*{Chemotherapy Extension}

The chemotherapy extension in Section~4 introduces a standard exponential
drug-decay equation and augments the majority of cellular compartments with a drug-kill
term. These terms where added under the assumption that chemotherapy attacks cancer and healthy cells alike. 
The assumption that chemotherapy does not
affect monocyte production is a simplification. 
Work by Larinova et al. \parencite{larionova2019interaction} 
summarizes the various effects of chemotherapy agents on macrophages in-vitro. 
Model expansion to explore the various dynamics is another avenue for further investigation. 
\\
\\
Two biologically distinct tumor-free scenarios were analyzed. The first
($T^* = 0$) corresponds to a post-treatment state in which cancer and drug have
been cleared but no residual immune activity remains. The
stability criterion for this case mirrors that of the macrophage extension,
with $\xi_2 = \beta M_n^* + \phi M_c^* - \delta$ again being the sole
potentially destabilizing eigenvalue. The second scenario ($T^* > 0$) represents
immune-mediated tumor
clearance with residual T-cell activity. Stability here is characterized by the
Routh-Hurwitz criterion $a_1 a_2 > a_3$ as described in Mathematical Biology by J.D. Murray \parencite{murray2002mathematical}.

\section{Conclusion}

This paper examined the tumor-immune model of Hernandez-Lopez et
al.\ \parencite{ERX2026}, identifying vague parameter description proposed biologically motivated corrections. Two
model extensions were then developed and analysed. The macrophage extension
captures the monocyte-to-macrophage differentiation  yielding a
four-variable system whose tumor-free equilibrium is characterized by the
combined immunostimulatory pressure $\beta M_n^* + \phi M_c^*$ relative to
lymphocyte death $\delta$. The chemotherapy extension further augments this
system with a drug concentration variable and analyses both immune-absent and
immune-active tumor-free states, demonstrating that post-treatment immune
activity provides an additional layer of stability governed by the
Routh-Hurwitz condition $a_1 a_2 > a_3$.
\\
\\
Collectively, these results indicate that model structure and parameter
interpretation impact the qualitative dynamics and
biological conclusions drawn from tumor-immune systems. In particular, the
sign of the monocyte--T-cell interaction parameter $\beta$ qualitatively
changes the nature of the tumor-free equilibrium, and the explicit inclusion
of macrophage dynamics reveals stability conditions that are invisible in
the reduced monocyte-only model. Future work would include bifurcation analysis 
using mathematical continuation software and simulation of biological viable parameters 
from experimental data, which due space and time limitations of this project the proof of 
concept code was not included but are available at https://github.com/Rosiedefazio/AMATH-881-Final-Project.

\newpage
\printbibliography

\end{document}